% Packages and macros needed by the chapters under PHIL/
% (extracted from the original paper preamble, PHIL/appendix.tex)

% T1 encoding: accented letters (ö in "Teichteil-Königsbuch") become single
% glyphs, so TeX can hyphenate the word instead of overflowing the margin.
\usepackage[T1]{fontenc}

\usepackage{booktabs}           % professional-quality tables
\usepackage{multirow}           % tabular cells spanning multiple rows
\usepackage{multicol}
\usepackage{amsfonts}           % blackboard math symbols
\usepackage{array}
\usepackage{wrapfig}
\usepackage{nicefrac}           % compact symbols for 1/2, etc.
\usepackage{url}
\usepackage{enumitem}
\usepackage{xspace}
\usepackage[mathscr]{euscript}
\usepackage[export]{adjustbox}

% Author-annotation helpers
\newcommand{\rex}[1]{{{\textcolor{magenta}{[rex: #1]}}}}
\newcommand{\jure}[1]{{{\textcolor{red}{[Jure: #1]}}}}
\newcommand{\michal}[1]{{{\textcolor{blue}{[Michal: #1]}}}}
\newcommand{\update}[1]{{{\textcolor{ForestGreen}{[Update: #1]}}}}

% Shorthands
\newcommand{\etal}{\textit{et al.}\xspace}
\newcommand{\alg}{\textsc{PHIL}\xspace}
\newcommand{\sail}[0]{\textsc{SaIL}\xspace}
\newcommand{\VS}{\mathcal{V}}
\newcommand{\xhdr}[1]{{\noindent\bfseries #1}.}

% Math
\DeclareRobustCommand{\bbone}{\text{\usefont{U}{bbold}{m}{n}1}}
\DeclareMathOperator{\EX}{\mathbb{E}} % expected value

% algorithm2e configuration (package itself is loaded in Disseration.tex)
\SetKwComment{Comment}{$\triangleright$\ }{}
\newcommand{\algorithmStyle}[0]{\footnotesize}


%%%%%%%%%%%%%%%%%%%%%%%%%%%%%%%%%%%%%%%%%%%%%%%%%%%%%%%%%%%%%%%%%%%%%%%%%%%%%%%
% HyBRiD chapter (from the ICML 2024 paper preamble)
%%%%%%%%%%%%%%%%%%%%%%%%%%%%%%%%%%%%%%%%%%%%%%%%%%%%%%%%%%%%%%%%%%%%%%%%%%%%%%%

\usepackage{mathtools}
\usepackage{amsthm}
\usepackage{tabularx}
\usepackage{caption}
\usepackage{pifont}
\usepackage{bbm}                % \mathbbm
\usepackage{microtype}
\usepackage[textsize=tiny]{todonotes}
\usepackage[capitalize,noabbrev]{cleveref}   % load last

% Theorem environments (numbered per chapter, not per section, in a book class)
\theoremstyle{plain}
\newtheorem{theorem}{Theorem}[chapter]
\newtheorem{proposition}[theorem]{Proposition}
\newtheorem{lemma}[theorem]{Lemma}
\newtheorem{corollary}[theorem]{Corollary}
\theoremstyle{definition}
\newtheorem{definition}[theorem]{Definition}
\newtheorem{assumption}[theorem]{Assumption}
\theoremstyle{remark}
\newtheorem{remark}[theorem]{Remark}

% Table rules/struts used by the HyBRiD tables
\newcommand{\Tstrut}{\rule{0pt}{2.9ex}}         % "top" strut
\newcommand{\Bstrut}{\rule[-1.2ex]{0pt}{0pt}}   % "bottom" strut
\newcommand{\TBstrut}{\Tstrut\Bstrut}
\newcommand{\shline}{\addlinespace[3pt]\hline\addlinespace[3pt]}

% Names
\newcommand{\mname}{\textsc{HyBRiD}}
\newcommand{\multiheadname}{\textit{multi-head drop-bottleneck}}
\newcommand{\avgimp}{$12.1\%$}

% Review annotations (\rex is defined above and deliberately not redefined)
\newcommand{\htl}[1]{\textcolor{orange}{[htl: #1]}}
\newcommand{\yize}[1]{\textcolor{magenta}{[Yize: #1]}}
\newcommand{\jialin}[1]{\textcolor{brown}{[jialin: #1]}}
\newcommand{\sw}[1]{{\color{magenta}***SW: #1***}}
\newcommand{\cam}[1]{{\color{red}#1}}


%%%%%%%%%%%%%%%%%%%%%%%%%%%%%%%%%%%%%%%%%%%%%%%%%%%%%%%%%%%%%%%%%%%%%%%%%%%%%%%
% MindLLM chapter (from the ICML 2025 paper preamble)
%%%%%%%%%%%%%%%%%%%%%%%%%%%%%%%%%%%%%%%%%%%%%%%%%%%%%%%%%%%%%%%%%%%%%%%%%%%%%%%
% Packages already loaded above are not repeated: mathtools, amsthm, tabularx,
% caption, pifont, cleveref, adjustbox, multicol, enumitem, array, subcaption.
% The paper's `subfigure` package is deliberately NOT loaded — it clashes with
% subcaption, and only the subcaption `subfigure` *environment* is used.
% Theorem environments are shared with the HyBRiD block above.

\usepackage{makecell}
\usepackage{siunitx}
\usepackage[normalem]{ulem}     % \uline
\usepackage[most]{tcolorbox}

% Names and marks
\newcommand{\mindllm}{MindLLM}  % the paper's \name, renamed to avoid a clash
\newcommand{\cmark}{\ding{51}}
\newcommand{\xmark}{\ding{55}}
\newcommand{\tdag}{\textsuperscript{\textdagger}}

% Review annotations (\rex is defined above and deliberately not redefined)
\newcommand{\yy}[1]{\textcolor{orange}{(Yujun: #1)}}
\newcommand{\qwk}[1]{\textcolor{blue}{(Weikang: #1)}}


%%%%%%%%%%%%%%%%%%%%%%%%%%%%%%%%%%%%%%%%%%%%%%%%%%%%%%%%%%%%%%%%%%%%%%%%%%%%%%%
% PRISM chapter (from the ICLR 2026 paper preamble)
%%%%%%%%%%%%%%%%%%%%%%%%%%%%%%%%%%%%%%%%%%%%%%%%%%%%%%%%%%%%%%%%%%%%%%%%%%%%%%%
% The paper's math_commands.tex defines 431 macros, none of which the chapters
% use, so it is not merged. Its `algorithm`/`algorithmic` packages are also
% skipped: no algorithm environment is used, and they would clash with the
% algorithm2e loaded for the PHIL chapter.

\newcommand{\modelname}{\textbf{{PRISM}}\xspace}


%%%%%%%%%%%%%%%%%%%%%%%%%%%%%%%%%%%%%%%%%%%%%%%%%%%%%%%%%%%%%%%%%%%%%%%%%%%%%%%
% RepQA chapter (from the npj Digital Medicine paper preamble)
%%%%%%%%%%%%%%%%%%%%%%%%%%%%%%%%%%%%%%%%%%%%%%%%%%%%%%%%%%%%%%%%%%%%%%%%%%%%%%%
% \bench is the only macro from that preamble the chapter actually uses; the
% rest (\model, \Set, \Mat, \vpara, \reb, \xm, \xp, prop/problem theorems) and
% its heavy packages (commath, listings, subfigure, threeparttable, lineno) go
% unused, so they are not merged. \rex/\htl already exist above.
\newcommand{\bench}{\textsc{RepQA}\xspace}

% A part-divider page that carries a short introductory paragraph on the SAME
% page as the part title (the book class would otherwise push the intro onto
% the following page). Usage: \partwithintro{Part title}{Intro paragraph...}
\makeatletter
\newcommand{\partwithintro}[2]{%
  \begingroup
    % Suppress the automatic \vfil\newpage that book.cls fires right after the
    % part title, so the introduction can share the part page.
    \renewcommand\@endpart{}%
    \part{#1}%
  \endgroup
  \vskip 2.5em
  \begin{quote}#2\end{quote}%
  \vfil            % balances the \null\vfil book.cls put before the title,
  \newpage         % centering the title+intro block vertically on the page.
}
\makeatother

% Last-resort stretch for paragraphs TeX cannot otherwise break — these papers
% were set in two columns and some have long unbreakable inline math that now
% overflows the wider one-column measure. Only applies where a line would
% otherwise stick out; well-set paragraphs are untouched.
\emergencystretch=3em
